\begin{abstract}
Video human mesh refiners routinely condition on 2D keypoints, yet the interface through which this evidence reaches the body is rarely isolated: keypoints are pooled or concatenated into one vector that is broadcast to every joint.
We ask whether preserving the identity of each landmark and routing it to the joints it constrains matters once the evidence is real.
An audit of our own joint-indexed feedback motivates the question: its landmark observations came from a head that never received gradient, so it routed drift from a frozen reference model rather than image evidence.
\method therefore turns external detector keypoints into identity-bearing tokens that carry the residual to the current body projection, the projection itself and the detector confidence.
Twenty-four SMPL joint queries read these tokens through attention with a kinematic-tree distance bias; the feedback is recomputed at every refinement iteration, enters before part pooling and masked temporal exchange, and starts from a zero-initialised output that preserves the inherited prediction.
On 3DPW, all variants start from one parent, share a 1,000-step budget and three seeds, and are tested once through validation locks; we compare routing with no feedback, capacity-matched pooling, ordered-vector conditioning, identity-free sets, unbiased and wrongly biased routing, and a weaker detector, using paired intervals that combine seed and test-sequence variance.
\TODO{One sentence with the main test result, \eg ``Routing lowers test MPJPE by \cmpdiff{B4-ext-joint}{B7-ext-concat}{test}{mpjpe}\,mm relative to ordered-vector conditioning (95\% CI \cmpci{B4-ext-joint}{B7-ext-concat}{test}{mpjpe})''; if the CIs include zero, state that identity-preserving routing performs on par with global conditioning and summarise what does matter.}
\end{abstract}
